\documentclass[11pt,a4paper]{article}

%========================
% PACKAGES
%========================
\usepackage[utf8]{inputenc}
\usepackage[T1]{fontenc}
\usepackage[french]{babel}
\usepackage[a4paper,margin=1.8cm]{geometry}
\usepackage{amsmath,amssymb,amsthm}
\usepackage{lmodern}
\usepackage{xcolor}
\usepackage{tikz}
\usepackage{pgfplots}
\usepackage{hyperref}
\usepackage{fancyhdr}
\usepackage{enumitem}

\pgfplotsset{compat=1.18}

%========================
% HEADER
%========================
\pagestyle{fancy}
\fancyhf{}
\fancyhead[L]{Mathématiques 3\ieme}
\fancyhead[C]{Chapitre : Nombres et Calculs}
\fancyhead[R]{\thepage}

%========================
% ENVIRONNEMENTS
%========================
\newtheorem{definition}{Définition}
\newtheorem{property}{Propriété}
\newtheorem{theorem}{Théorème}
\newtheorem{example}{Exemple}
\newtheorem{remark}{Remarque}

\title{\Huge Nombres et Calculs \\ \large Cours complet niveau 3\ieme}
\date{}

\begin{document}

\maketitle

\section{Pourquoi ce chapitre est fondamental ?}

Ce chapitre est le \textbf{socle des mathématiques}.  
Tout ce que tu feras ensuite (fonctions, physique, informatique…) repose dessus.

\textbf{Idée clé :}  
\begin{center}
\textit{Les mathématiques sont un langage pour décrire le réel.}
\end{center}

Exemples :
\begin{itemize}
\item partager une pizza → fractions
\item croissance d’un virus → puissances
\item calculer un prix → équations
\end{itemize}

---

# =========================
\section{Calcul numérique}

\subsection{Fractions}

\begin{definition}
Une fraction est un quotient :
\[
\frac{a}{b}, \quad b \neq 0
\]
\end{definition}

\subsubsection{Interprétation visuelle}

\begin{center}
\begin{tikzpicture}
\draw (0,0) circle (2cm);
\draw (0,0) -- (2,0);
\draw (0,0) -- (0,2);
\draw (0,0) -- (-2,0);
\draw (0,0) -- (0,-2);
\fill[blue!30] (0,0) -- (2,0) arc (0:90:2) -- cycle;
\end{tikzpicture}
\end{center}

\textbf{Analogie :}  
$\frac{1}{4}$ = une part d’une pizza coupée en 4.

---

\subsection{Simplification}

\begin{property}
\[
\frac{a}{b} = \frac{ka}{kb}
\]
\end{property}

\textbf{Démonstration :}
\[
\frac{ka}{kb} = \frac{k}{k} \times \frac{a}{b} = \frac{a}{b}
\]

---

\subsection{Puissances}

\begin{definition}
\[
a^n = a \times a \times \cdots \times a
\]
\end{definition}

\textbf{Exemple réel :}

Si tu doubles ton argent chaque jour :
\[
2^n
\]

Jour 10 → $2^{10} = 1024$

�� croissance \textbf{explosive}

---

\subsection{Racines carrées}

\begin{definition}
\[
\sqrt{a} = b \Leftrightarrow b^2 = a
\]
\end{definition}

\textbf{Interprétation géométrique :}

\begin{center}
\begin{tikzpicture}
\draw (0,0) rectangle (3,3);
\node at (1.5,1.5) {$9$};
\end{tikzpicture}
\end{center}

Aire = 9 → côté = $\sqrt{9} = 3$

---

# =========================
\section{Calcul littéral}

\subsection{Développement}

\begin{property}
\[
a(b+c) = ab + ac
\]
\end{property}

\textbf{Démonstration géométrique :}

\begin{center}
\begin{tikzpicture}
\draw (0,0) rectangle (6,2);
\draw (3,0) -- (3,2);
\node at (1.5,1) {$ab$};
\node at (4.5,1) {$ac$};
\end{tikzpicture}
\end{center}

---

\subsection{Identités remarquables}

\begin{property}
\[
(a+b)^2 = a^2 + 2ab + b^2
\]
\end{property}

\textbf{Démonstration visuelle :}

\begin{center}
\begin{tikzpicture}
\draw (0,0) rectangle (4,4);
\draw (0,3) -- (4,3);
\draw (3,0) -- (3,4);
\end{tikzpicture}
\end{center}

---

\subsection{Factorisation}

\begin{definition}
Factoriser = écrire sous forme de produit.
\end{definition}

\begin{example}
\[
3x + 6 = 3(x+2)
\]
\end{example}

\textbf{Analogie :}  
mettre des objets dans une boîte commune.

---

# =========================
\section{Équations}

\subsection{Principe fondamental}

\begin{center}
\textbf{On fait la même chose des deux côtés}
\end{center}

\textbf{Analogie : balance}

\begin{center}
\begin{tikzpicture}
\draw (-2,0) -- (2,0);
\draw (0,0) -- (0,-1);
\draw (-1,-1) -- (-1,-2);
\draw (1,-1) -- (1,-2);
\end{tikzpicture}
\end{center}

---

\subsection{Résolution}

\begin{example}
\[
2x + 3 = 7
\]
\[
2x = 4
\Rightarrow x = 2
\]
\end{example}

---

# =========================
\section{Inéquations}

\begin{property}
Si on multiplie par un nombre négatif :
\[
\text{on inverse le sens}
\]
\end{property}

\begin{example}
\[
-2x > 4 \Rightarrow x < -2
\]
\end{example}

---

# =========================
\section{Arithmétique}

\subsection{Diviseurs}

\begin{definition}
$a$ divise $b$ si :
\[
b = ak
\]
\end{definition}

---

\subsection{Nombres premiers}

\begin{definition}
Un nombre premier possède exactement 2 diviseurs.
\end{definition}

\textbf{Liste :}
\[
2,3,5,7,11,13,17
\]

\textbf{Analogie :}
Ce sont les \textbf{atomes des mathématiques}.

---

\subsection{Décomposition en facteurs premiers}

\begin{example}
\[
60 = 2^2 \times 3 \times 5
\]
\end{example}

\textbf{Clé :}
Comme décomposer une molécule en atomes.

---

# =========================
\section{Applications concrètes}

\begin{itemize}
\item économie : calcul de prix
\item physique : formules
\item informatique : algorithmes
\item ingénierie : modélisation
\end{itemize}

---

# =========================
\section{Conclusion}

\textbf{À maîtriser absolument :}
\begin{itemize}
\item fractions
\item puissances
\item développement / factorisation
\item équations
\item nombres premiers
\end{itemize}

\bigskip

\end{document}